%% notes-tex/024-perturb-seq-pilot/024-perturb-seq-pilot.tex
%% [[experiments.024-perturb-seq-costing.pilot-plan]]
%% https://github.com/Mjvolk3/torchcell/tree/main/notes-tex/024-perturb-seq-pilot/024-perturb-seq-pilot.tex
%%
%% The bench plan that falls out of notes-tex/024-perturb-seq-costing. That
%% document prices a genome-scale screen and ranks what is unknown; this one
%% names the two runs to start, one per host, and answers the two questions the
%% costing does not: whether to buy environments or a readout first, and whether
%% to raise plasmid copy number early.
%%
%% Written for the people who will run it. Every number is cross-referenced to
%% the section or script that derives it, and nothing new is asserted here.
%%
%% Build:  make            -> 024-perturb-seq-pilot.pdf       (draft view)
%%         make clean-view -> 024-perturb-seq-pilot-clean.pdf (share view)
%%         make check      -> formatting / provenance / citation gate

\documentclass[10pt,a4paper]{article}
\usepackage[draft]{tcdoc}

\title{\vspace{-1.5em}First Perturb-seq Experiments in Two Yeasts\\
  {\large A bench plan for \textit{S. cerevisiae} and \textit{P. kudriavzevii}}}
\author{Michael Volk}
\date{\today}

\begin{document}
\maketitle
\thispagestyle{empty}

\begin{abstract}
\noindent
Two first runs are ready to start, one per host, and they do not depend on each
other. In \org{S.\ cerevisiae} the guide library already exists, so the open
question is whether perturbation identity can be read out of a cell at all, and
one droplet channel answers it for \$4{,}493. In \org{P.\ kudriavzevii} no
pooled library is reachable yet, so the first run carries no guides and instead
measures whether the assay works on that cell wall, which is worth doing across
several growth media rather than one. Both are sized against a published
genome-scale yeast screen whose data, on reanalysis, does not reproduce itself,
and the three checks that would have caught that are built into the designs
here.
\end{abstract}

\vspace{0.5em}
{\small\color{tcgray}
\textbf{How to read this.} \Cref{sec:plan} is the plan and the costs.
\Cref{sec:routes} is what the bench actually does on each of the two routes.
\Cref{sec:nadal} is the published screen and what went wrong with it, which is
where the quality checks in \cref{sec:checkpoints} come from. \Cref{sec:open}
answers the two decisions that were open: media against readout, and plasmid
copy number.\par\smallskip

Costing and method review: \file{notes-tex/024-perturb-seq-costing/}\par
Generating scripts: \file{experiments/024-perturb-seq-costing/scripts/pilot\_options.py}\par
References: Zotero \file{torchcell/microbe-perturb-seq} (group $\cup$ personal)\par}

\vspace{1em}
\tccontents

%% notes-tex/024-perturb-seq-pilot/sections/1-plan.tex
%% SOURCE: experiments/024-perturb-seq-costing/scripts/pilot_options.py
%%         (results/pilot_options.json; tables/t1-pilot-options.tex is emitted by it)
%% Every dollar figure composes cost_model.py's platforms over the Carver Center
%% rate card in uiuc_core_data.py. No price is entered here by hand.

\section{The Plan\secstatus{todo}}
\label{sec:plan}

Two runs, one per host, started in parallel by two people. Neither waits on the
other, and they answer different questions because the two organisms are at
different stages.

\notesec{\org{S.\ cerevisiae}: the library exists, so measure the readout} A
genome-wide single-guide CRISPR interference (\term{CRISPRi}) library is already
in hand, six guides per gene over roughly 6{,}000
genes~\citep{lianMultifunctionalGenomewideCRISPR2019}. What has never been
established is whether a cell's guide can be recovered from its RNA. Call that
probability \term{q}, the fraction of cells whose perturbation is read. One
channel on the core's 10x Genomics 5$'$ kit answers it (\cref{sec:guide}).

\notesec{\org{P.\ kudriavzevii}: no library is reachable, so measure the assay}
Transformation efficiency in this host is three to five orders of magnitude
below \org{S.\ cerevisiae}, which puts a pooled library-scale screen out of
reach until that moves. A first run therefore carries no guides at all, and what
it returns is whether the cell wall yields to in-droplet digestion, how many
messenger RNA molecules are recovered per cell, and what fraction of reads is
ribosomal. None of the three needs genetics.

%% GENERATED FILE -- do not hand-edit.
%% SOURCE: experiments/024-perturb-seq-costing/scripts/pilot_options.py
\begin{table}[htbp]\centering
\footnotesize
\caption[]{\textbf{The candidate first runs, priced end to end.} Usable cells are those expected to survive quality control, 55\% of a droplet channel's 20{,}000 and 25\% of a protocol run's 480{,}000. Recurring cost is reagents plus sequencing at the cheapest configuration the Carver Center sells that holds the reads, which at this scale is never the flow cell a genome-scale screen would use. The one-time column is the split-pool barcode plate set, which covers about 215 protocol runs. The route is the plate one for the eight-media row and 10x droplet for every other.}\label{tab:pilots}
\begin{tabular}{@{}>{\raggedright\arraybackslash}p{40mm}>{\raggedright\arraybackslash}p{23mm}rrr>{\raggedright\arraybackslash}p{44mm}@{}}
\toprule
run & host & usable cells & recurring & one-time & returns \\
\midrule
Guide capture, 1 channel & \org{S. cerevisiae} & 11,000 & \$4,493 & -- & q, the per-cell guide detection rate \\
Guide capture, low vs high copy & \org{S. cerevisiae} & 22,000 & \$8,576 & -- & q on two vector backbones \\
Unperturbed profile, 1 medium & \org{P. kudriavzevii} & 11,000 & \$3,810 & -- & wall digestion, UMIs per cell, rRNA fraction \\
Unperturbed profile, 4 media & \org{P. kudriavzevii} & 44,000 & \$12,550 & -- & the same three, plus a 4-condition response atlas \\
Plate route, 8 media in one run & \org{P. kudriavzevii} & 20,000 & \$4,990 & \$10,000 & the same three, across 8 media, plus the plate chemistry \\
\bottomrule
\end{tabular}
\end{table}


%% SOURCE: experiments/024-perturb-seq-costing/scripts/pilot_options.py
\tcfigfit[1.0]{figures/pilot_options.pdf}{%
  \textbf{The two first runs cost under \$9{,}000 together, and the cheaper way
  to add a growth medium is a plate rather than a second droplet channel.}
  \textbf{(a)} Recurring cost of each candidate, split into reagents (solid, by
  host) and sequencing (hatched). Sequencing is priced on the cheapest
  configuration the Carver Center sells that holds the reads, which at this
  scale is a MiSeq run or a split NovaSeq lane rather than the large flow cell a
  genome-scale screen fills. The plate route carries a \$10{,}000 one-time
  purchase of barcoded oligo plates, which covers about 215 protocol runs.
  \textbf{(b)} Recurring cost against the number of growth media, at one droplet
  channel's worth of cells per medium. An unmodified droplet channel carries
  nothing that says which medium a cell came from, so each medium buys its own
  channel, at \$2{,}650 for the first and \$2{,}240 for each of the next six.
  Both plate routes read a
  well index out of every cell already, so media ride together in one run and
  the steps are sublibrary and run rounding rather than per-medium cost.
  \textbf{(c)} What one run buys against what a screen needs, on usable cells.
  \textbf{(d)} Why $q$ is the quantity worth measuring first. At one guide per
  cell it scales the yield linearly; at $k$ guides per cell the usable fraction
  is $q^{k}$, so the two published microbial values, 0.21 in
  \org{E.\ coli}~\citep{brandnerPooledSinglecellCRISPRa2025} and more than 0.71
  in \org{S.\ cerevisiae}~\citep{nadal-ribellesSinglecellResolvedGenotypephenotype2025},
  are the difference between multiplexing being expensive and being impossible.
  Generated by \file{experiments/024-perturb-seq-costing/scripts/pilot\_options.py}.}{fig:pilots}

\notesec{What the prices do and do not include} Reagents and sequencing at
on-campus rates, effective 2026-08-01. Not included: strain construction, the
media themselves, staff time, and the cell counter discussed in
\cref{sec:counting}. The sequencing line is the one most likely to move, because
a lane is an indivisible purchase and a run that half-fills one pays for the
whole thing.

\begin{keybox}
\textbf{Start both. They cost \$8{,}303 together and answer different
questions.} In \org{S.\ cerevisiae}, one 5$'$ channel on the existing CRISPRi
library measures $q$, the fraction of cells whose guide is readable, at
\$4{,}493. In \org{P.\ kudriavzevii}, one channel with no library measures wall
digestion, messenger RNA per cell, and ribosomal fraction, at \$3{,}810. Neither
depends on the other, one person can own each, and the two results together
decide every design question after them.
\end{keybox}

%% notes-tex/024-perturb-seq-pilot/sections/2-routes.tex
%% SOURCE: hand-authored synthesis of notes-tex/024-perturb-seq-costing/
%%         sections/3-method-explainers.tex (Secs. 3.1, 3.2, 3.7) and
%%         sections/5-application.tex (Sec. 5.5, the core meeting). No new
%%         numbers; each is cross-referenced to the section that derives it.

\section{The Three Bench Routes\secstatus{todo}}
\label{sec:routes}

All three routes end in the same object, a table of cells by genes with a
perturbation label on each cell. They differ in how a cell keeps its identity
while its RNA is being converted to a sequencing library, and that difference
decides who does the work. Two of them are the familiar pair, plate barcoding
and droplet isolation. The third composes them, and it is the one this program
is aiming at.

%% SOURCE: notes/assets/drawio/perturb-seq-workflow.drawio.svg, exported by
%% `make figures`. Panel numbering and content are unchanged from Fig. 5 of
%% 024-perturb-seq-costing.
\tcfigfit[1.0]{figures/perturb-seq-workflow.pdf}{%
  \textbf{The whole experiment, from library to count matrix, and the one stage
  where the three routes differ.} Read left to right. Two symbols recur in panel
  (b): \textbf{H} is the PCR handle, attached at one end only during barcoding,
  and \textbf{CB} is the cell barcode. The color key is repeated at the foot of
  the figure.
  \textbf{(a)} The stages. Every wet-lab step is common to all three routes
  except the shaded block, where split-pool barcodes cells across three rounds of
  plates, unmodified droplet gives each cell its own bead, and preindexing does
  one plate round first and then a droplet, so several cells can share a droplet
  and still be told apart. Fragmentation is drawn as its own stage for two
  reasons: its position relative to ribosomal RNA depletion decides how much
  that depletion is worth, roughly \$99{,}000 at genome scale, and it determines
  what a count means.
  \textbf{(b)} The same physical cutting, with opposite consequences, which is
  the reason this figure is worth reading rather than skimming. In bulk RNA
  sequencing every piece of a molecule can be sequenced, so a long transcript
  yields proportionally more reads and the length has to be divided back out,
  which is what FPKM and TPM exist to do. In a 3$'$-counting single-cell assay
  only the barcode-proximal piece keeps the outer handle, so exactly one
  fragment survives per molecule whatever its length. Applying TPM or FPKM to a
  count matrix from either route is therefore an error rather than a stylistic
  choice.
  \textbf{(c)} The funnel, which is what a budget is actually made of. Cells are
  lost at recovery, at quality control and at guide assignment, and reads are
  lost separately to the spike-in, to ribosomal sequence and to cells that are
  discarded later. Composed from the protocols rather than traced from a
  published figure, and shared with the costing review, where the numbers drawn
  in it are sourced individually.}{fig:workflow}
%% Source notes/assets/drawio/perturb-seq-workflow.drawio, exported by `make figures`

\notesec{Droplet, which the core runs as a service} A microfluidic chip
co-encapsulates one cell with one barcoded bead, and every molecule released
inside that droplet picks up the bead's barcode. The Carver Center owns the
instrument, the GEM-X 5$'$ kit and the quality control, so our bench work
ends when fixed cells are delivered. Yeast needs one modification, roughly
a microlitre of zymolyase in the reaction to digest the cell
wall~\citep{jarianiNewProtocolSinglecell2020}, and the core has said it accepts
that. A channel recovers about 20{,}000 cells and costs \$2{,}650 for the first
library and \$2{,}240 for each of the next six.

\notesec{Plate, which we would run ourselves} Cells are fixed and permeabilized,
then distributed across a 96-well plate where each well appends a different
barcode; the wells are pooled, re-split and barcoded again, twice more. A cell's
identity is the sequence of wells it visited, so no instrument ever has to
isolate it~\citep{brettnerUltraHighthroughputMassively2024}. One pass through
the protocol handles up to 96 samples and, at a full plate, roughly 480{,}000
cells. The reagents are cheap and the labor is not: about a week of skilled
bench work per run, against a channel the core runs.

\notesec{Preindexed droplet, which is one plate round and then a droplet}
\label{sec:preindexed}
The third route composes the other two, and it is the one worth aiming at.
Permeabilized cells are barcoded once on a 384-well plate, where a well-specific
primer reverse transcribes inside the still-intact cell and writes round one into
the complementary DNA. The cells are then pooled and deliberately overloaded into
a single droplet channel, where a bead writes round
two~\citep{datlingerUltrahighthroughputSinglecellRNA2021}. Neither round
identifies a cell on its own, and that is the entire trick: a droplet may hold
several cells, because the plate already told them apart.

What that removes is the constraint that makes droplet screening expensive. An
ordinary channel runs at low occupancy so that two cells rarely share a droplet,
which leaves most droplets empty and most of the reagent unused. Datlinger et
al.\ loaded 100-fold over the recommendation, filled 95.5\% of droplets at an
average of 9.6 nuclei each, and recovered \textbf{151{,}788} transcriptomes from
one channel, which is 7.6 times the 20{,}000 the core quotes. Since reagent is
sold per channel, a genome-scale screen falls from 137 channels and \$364{,}000
to 18 channels and \$109{,}000. Depleted split-pool is still cheaper at
\$57{,}500, but a factor of 1.9 is small enough to be decided by things the
dollar column does not hold, and all of them favor this route: the core runs the
droplet platform as a service and owns its quality control, the 384-well plate
labels growth conditions for nothing, and an accurate cell count matters far less
when the channel is saturated on purpose.

\textbf{The catch is that it has never been run in any microbe.} Every
demonstration is human or mouse. It also inherits the plate route's sample
preparation rather than the droplet route's: writing a barcode inside a cell
needs reagents to reach the transcriptome in situ, which in yeast means fixation
and wall digestion rather than the in-droplet zymolyase above.
\textbf{Hypothesis (untested):} the two halves compose, since in-cell reverse
transcription on fixed yeast and droplet capture of yeast are each established
separately.

\notesec{A first barcode round on a plate is what makes growth media cheap}
Wherever round one is a plate of wells, its index is read out of every cell
whether or not it has been given a meaning. Growing conditions in a deep-well
plate and barcoding them into the same run therefore costs no barcode space at
all (\cref{fig:pilots}b), and this is the property the two plate-first routes
share: split-pool carries 96 such labels and preindexing carries 384. Only the
unmodified droplet route lacks them, so there each condition buys its own
channel. That is what makes a plate-first route the one to use when the
experiment is about growth media, and it is why the \org{P.\ kudriavzevii} run
should not be a single condition.

\notesec{The guide readout, which is the whole \org{S.\ cerevisiae} question}
\label{sec:guide}
Guide RNAs are transcribed by RNA polymerase III, so they are neither capped nor
polyadenylated, and an assay that captures messenger RNA by its poly(A) tail
sees nothing. That is why our library is read today by amplifying the plasmid
rather than the RNA. The 5$'$ kit may not have the problem: it captures the
guide with a primer against the guide's own scaffold, carried in the
reverse-transcription mix, and 10x state the requirement as an architectural one,
that guides be ``engineered for use with standard Cas9 systems with a
protospacer on the 5$'$ end''\external{10x CG000735}. A conventional polymerase
III cassette, including ours, has exactly that architecture.

\textbf{Hypothesis (untested):} the scaffold primer reverse transcribes a yeast
polymerase III guide, and perturbation identity is therefore readable from an
unmodified library. Two things have to hold and neither is established. The
primer set is proprietary and validated against the standard \gene{SpCas9}
scaffold, and 10x say plainly that ``alternate sgRNA structures \dots\ have not
been tested''\external{10x CG000735}, so our scaffold has to match theirs. And
the poly-dT half of that primer mix contributes nothing to a guide with no tail,
so capture rests entirely on the scaffold primer. If it fails, the fallback is a
polyadenylated guide barcode transcribed by polymerase II, which is a
construction project of several months.

The guide library is sequenced far more shallowly than the transcriptome, 5{,}000
read pairs per cell against 20{,}000, and 10x warn against sequencing it
alone\external{10x CG000735}, so it is pooled with the expression library rather
than submitted separately. It is in the price in \cref{tab:pilots}.

\begin{keybox}
\textbf{One channel, one existing library, and the answer is a number.} The
\org{S.\ cerevisiae} run is worth starting before anything else because it is
the cheapest way to find out whether the whole program needs a new construct.
If $q$ comes back near the 0.71 a published yeast screen reports, single-guide
screening works today and multiplexing is merely expensive. If it comes back
near zero, the 5$'$ scaffold primer does not engage a polymerase III transcript
and a polyadenylated guide barcode has to be built before any screen is
designed.
\end{keybox}

%% notes-tex/024-perturb-seq-pilot/sections/3-nadal.tex
%%
%% SOURCE: experiments/028-knockout-expression/scripts/
%%           cross_study_ko_expression.py     -> results/cross_study_ko_expression.json
%%           cross_study_recomputed.py        -> results/cross_study_recomputed.json
%%           nadal_batch_replication.R        -> cross-batch and split-half tables
%%           nadal_replication_coverage.py    -> results/nadal_replication_coverage.json
%%           nadal_identify_deletion.py       -> results/nadal_identify_deletion.json
%%           nadal_paper_deleteome_comparison.R
%%           gene_covariation_all.py          -> results/gene_covariation_all.json
%%         gene_covariation_all.py is quoted in prose rather than figured, to
%%         hold the document to four figures; its panel is Fig. 5 of the 019
%%         document.
%%         Those scripts and their results live on branch
%%         `multimodal-phenotype-retrospective` (PR #265); the typeset treatment
%%         is notes-tex/019-simb-multimodal-expression, Sec. "A third deletion
%%         panel". The three figure PDFs here are that build's, committed
%%         unchanged. No number in this section is entered by hand.

\section{The Published Screen, and Why We Do Not Trust It\secstatus{todo}}
\label{sec:nadal}

The closest published experiment to the one proposed here is a genome-scale
single-cell screen of the yeast knockout collection, about 3{,}100 deletions
profiled by single-cell RNA
sequencing~\citep{nadal-ribellesSinglecellResolvedGenotypephenotype2025}. It is
the design this plan would otherwise copy. Loading it into our database and
comparing it against the knockout expression data already there showed it does
not agree with them, and then that it does not agree with itself.

This matters at the bench rather than in the analysis. Every failure below is
one a run can be designed to catch, and all three checks are cheap.

\notesec{First, what agreement looks like when it works} Two microarray studies
of the same deletion collection, measured years apart in different
hands~\citep{kemmerenLargeScaleGeneticPerturbations2014,sameithHighresolutionGeneExpression2015},
agree on a deletion's expression profile at a median correlation of \textbf{0.74}
over the 82 deletions they share, and on a reporter gene's behavior across
deletions at 0.62. That is the reference number: a real knockout signature
survives a change of laboratory.

\notesec{Against that, the single-cell panel reads as noise} Over the 914
deletions it shares with the microarrays, the median agreement is
\textbf{0.003}. Recomputing the fold change from the raw counts, rather than
taking the released values, moves it to 0.013. The agreement does not improve
with the number of cells behind a profile, which is what an undersampled but
real measurement would do.

%% SOURCE: experiments/028-knockout-expression/scripts/cross_study_ko_expression.py
\tcfigfit[0.92]{figures/cross_study_ko_expression.pdf}{%
  \textbf{A deletion's microarray profile reproduces in a second microarray study
  at a median correlation of 0.74 and in the single-cell panel at 0.00.}
  \textbf{(a, b)} The released single-cell values sit on a different scale, and
  one cell in five is a sentinel produced when a gene's mean is zero in one
  group, which the fold-change formula returns as roughly $\pm 23$. Those are
  discarded before every comparison here. \textbf{(c)} Spread of each deletion's
  profile. \textbf{(d)} Every shared deletion-by-gene value, with the fitted
  slope. \textbf{(e)} The headline: per-deletion agreement for the two
  microarrays against the microarray-to-single-cell comparison. \textbf{(f)} The
  same per reporter gene. \textbf{(g, h, i)} What the agreement tracks. It rises
  with how strongly a deletion responds and with how variable a reporter is, and
  it does not track cell count. Generated by
  \file{experiments/028-knockout-expression/scripts/cross\_study\_ko\_expression.py}.}{fig:crossstudy}

\notesec{The panel does not reproduce itself, which is the finding that
settles it} A cross-study zero cannot be interpreted without a within-study
replicate, and the published work reports none. One is recoverable, because the
same genotype's cells fall in several capture cartridges. Comparing a genotype
against itself across two cartridges, each referenced to wild-type cells from
its own cartridge, gives a median correlation of \textbf{0.043} over 17{,}554
pairs. Two \emph{different} genotypes compared the same way give 0.022. Splitting
one genotype's cells in half within a single cartridge gives 0.080 against a
0.051 null.

A genotype's profile reproduces itself about 0.02 above an unrelated genotype.
Nothing downstream of that can agree with anything.

Two related traps are worth naming because both are easy to repeat. Referencing
every profile to one pooled wild type instead of a per-batch wild type raises
both the real and the null value to 0.11, which looks like signal and is a
batch effect: wild type against wild type between cartridges differs by a
standard deviation of 0.75 in $\log_2$, with 16\% of genes beyond two-fold. And
sharing the reference between two halves of a split-half test inflates it to
0.37, which is the reference's own sampling noise measured twice.

%% SOURCE: experiments/028-knockout-expression/scripts/nadal_replication_coverage.py
%% and nadal_batch_replication.R (the cross-batch and split-half tables it reads)
\tcfigfit[0.92]{figures/nadal_replication_coverage.pdf}{%
  \textbf{The same genotype measured twice agrees no better than two different
  genotypes, and the published agreement statistic is a cell-count artifact.}
  \textbf{(a)} A genotype's profile in one capture cartridge against its profile
  in another, each referenced to wild-type cells from its own cartridge, with
  the different-genotype null beside it. The pooled-reference version of the
  same comparison is shown to its right, where a common batch effect lifts both.
  \textbf{(b)} The same within one cartridge, splitting the genotype's cells and
  the wild-type cells in half so the two profiles share no cell. \textbf{(c, d)}
  Coverage and depth: a median cell carries 934 unique molecules over 492 genes,
  and a median genotype has 92 cells. \textbf{(e)} The comparison the paper
  publishes, a count of changed genes against the microarray's count, which
  correlates at 0.23. \textbf{(f)} The profile correlation the paper's own
  released script computes and does not show, median 0.013. Its changed-gene
  count falls with cell number at a rank correlation of $-0.67$, because a
  genotype with few cells has more genes at a zero mean and each one becomes a
  sentinel that passes the fold-change threshold. Generated by
  \file{experiments/028-knockout-expression/scripts/nadal\_replication\_coverage.py}.}{fig:replication}

\notesec{The positive control fails too} A deletion removes its own transcript,
so in a clean profile the deleted gene sits at the bottom. In the microarray
data it is the single most reduced gene in \textbf{59.5\%} of 1{,}479 strains
and in the bottom one percent in 93\%. In the single-cell panel it is the most
reduced gene in under 1\% of genotypes, and its median rank sits near the middle
of the profile. Direct measurement of label purity puts roughly \textbf{40\%} of
the cells under a label as the genotype that label names.

\notesec{Everything else in our database agrees with everything else} Five
independent yeast panels, two knockout microarray studies, a knockout proteome,
a natural-isolate transcriptome survey and this screen, can be compared without
sharing a single strain, by asking whether the same pairs of genes move together
within each panel. The four established panels agree with one another at 0.30 to
0.68. This one agrees with all of them at 0.03 to 0.10, and its gene-pair
structure is nearly flat.

\begin{keybox}
\textbf{The lesson is three cheap checks, not an argument against single-cell
screening.} What went wrong is measurable from the data itself, so build the
measurement in. Replicate at least one genotype across two independent barcode
wells or cartridges and reference each to wild-type cells from the same one.
Read the targeted gene's own transcript as a positive control, since CRISPR
interference should reduce it in the cells carrying that guide. And report cells
per label and label purity beside every result. The screen above would have
failed all three before any biology was claimed from it.
\end{keybox}

%% notes-tex/024-perturb-seq-pilot/sections/4-open.tex
%% SOURCE: hand-authored synthesis. The environment arithmetic is
%%         experiments/024-perturb-seq-costing/scripts/scaling_analysis.py
%%         (Sec. 7.3 of the costing review) and pilot_options.py (Fig. 1b); the
%%         delivery facts are Sec. 7.1 of the same review. No new numbers.

\section{The Two Decisions That Were Open\secstatus{todo}}
\label{sec:open}

\subsection{Growth Media, or the Readout First\secstatus{todo}}
\label{sec:media}

Both, and the split falls along the two hosts rather than along a preference.

Growth media are the cheap axis. Cell demand grows in proportion to the number
of conditions, where covering named pairs of genes grows with the square of the
panel, and any route whose first barcode is a plate of wells labels conditions
for nothing (\cref{fig:pilots}b). A model that proposes strain designs also
needs both halves of its input to vary, and a dataset that varies only genotype
in one medium constrains only the genotype half.

What media cannot do is multiply an assay that does not exist yet. Running four
conditions through a readout whose guide detection is unmeasured produces four
copies of the same unanswered question, at four times the price. So the ordering
is not a judgment about which is more interesting.

\begin{itemize}[leftmargin=1.4em,itemsep=2pt,topsep=2pt]
  \item \textbf{\org{S.\ cerevisiae} spends its first run on the readout}, in
    one condition, because the library exists and $q$ is the gate
    (\cref{sec:guide}). Media come second, and by then the plate route is the
    way to buy them.
  \item \textbf{\org{P.\ kudriavzevii} spends its first run on media}, because it
    has no library to gate on. A single unperturbed condition under-uses a
    channel that has already been paid for, and the honest way to thicken that
    run is more conditions rather than genetics that do not exist yet.
\end{itemize}

On the plate route eight media ride in one protocol run for \$4{,}990 plus the
one-time plate purchase, against \$12{,}550 for four media as four droplet
channels (\cref{tab:pilots}). That comparison is the argument for the plate, and
it only applies once the plate chemistry is known to work on this cell wall,
which nobody has shown. Running the \org{P.\ kudriavzevii} first experiment on
droplet across a few media, rather than on the plate across many, buys the wall
answer on a platform the core supports and leaves the plate for run two.

\subsection{Plasmid Copy Number\secstatus{todo}}
\label{sec:copy}

Not early, with one exception worth building into the first run.

\notesec{Copy number is not guide diversity, and conflating them is the trap} A
cell that takes up one plasmid molecule and amplifies it to fifty copies carries
fifty copies of \emph{one} guide. Diversity comes only from the number of
independent uptake events at transformation. Raising average copy number by
weakening the selection marker is well established and finely tunable in yeast,
over ``5--100 copies per
cell''~\citep{lianConstructionPlasmidsTunable2016}, and it does not by itself
put two guides in a cell.

\notesec{Three reasons to keep the pilot low-copy} A single-guide screen is
interpretable precisely because one cell means one guide, and the field builds
its libraries in ``typically, low-copy number vectors to prevent multiple guide
plasmids in a single
cell''~\citep{robertsonAdvancesCRISPRenabledGenomewide2025}. Multiple plasmids
per cell are also transient: multi-vector transformants segregate over roughly a
day of outgrowth~\citep{scanlonQuantifyingResolvingMultiple2009}, and 2$\mu$
plasmids are retained in only about 60\% of cells even in a laboratory
strain~\citep{haysNaturalVariantEssential2020}, so a high-copy pilot measures a
population that is changing while it is being measured. And a cell carrying two
guides at unknown copy number confounds the one number the first run exists to
produce, since a guide missed in such a cell and a guide missed in a
single-guide cell are not the same failure.

\notesec{The exception, and it is a hypothesis} Guide detection depends on how
much guide transcript is present, and copy number raises that.
\textbf{Hypothesis (untested):} a higher-copy vector raises $q$ by putting more
guide RNA in each cell. Nothing measures this, and it is directly testable in
the run already planned, as a second channel carrying the same library on a
higher-copy backbone (\cref{tab:pilots}, row two, \$8{,}576 for the pair rather
than \$4{,}493). Worth doing if the budget allows, and worth skipping if the
first channel returns a high $q$ on its own.

\notesec{When copy number does become the right lever} Once $q$ is known and
single-guide screening works, putting several guides in one cell is what makes
combinations affordable, and the multi-plasmid route is the only way past the
two-to-four guide ceiling that cloning arrays into one vector imposes. The
missing measurement then is the per-cell distribution of distinct guides under
tuned delivery, which the literature does not report for yeast. That is the
second pilot, not the first.

\begin{keybox}
\textbf{Keep the first library low-copy, and treat high copy as a detection
experiment rather than a multiplexing one.} Copy number buys guide transcript
per cell, not guides per cell, so it cannot deliver combinations and it can
confound the measurement of $q$. The one version worth running now is a second
channel of the same library on a higher-copy backbone, purely to see whether $q$
rises with it.
\end{keybox}

%% notes-tex/024-perturb-seq-pilot/sections/5-checkpoints.tex
%% SOURCE: hand-authored synthesis. Sequencing prices and cell counts are from
%%         experiments/024-perturb-seq-costing/scripts/pilot_options.py and
%%         uiuc_core_data.py; the counting and core-service facts are Sec. 5.5 of
%%         the costing review; the three checks are Sec. 3 of this document.

\section{Running It\secstatus{todo}}
\label{sec:checkpoints}

\subsection{What Each Run Must Return\secstatus{todo}}
\label{sec:returns}

A run is finished when it has produced these, and not when it has produced data.

\notesec{\org{S.\ cerevisiae}, one 5$'$ channel on the existing CRISPRi library}
\begin{enumerate}[leftmargin=1.6em,itemsep=2pt,topsep=2pt]
  \item $q$, the fraction of quality-controlled cells carrying exactly one
    identified guide. This is the number the run exists for.
  \item The targeted gene's own transcript in the cells carrying its guide,
    against cells carrying other guides. Interference should reduce it, and this
    is the positive control whose absence went unnoticed in \cref{sec:nadal}.
  \item Messenger RNA molecules and genes per cell, and the fraction of reads
    that are ribosomal.
  \item Guide representation: how many of the library's guides appear at all, and
    how unevenly.
\end{enumerate}

\notesec{\org{P.\ kudriavzevii}, one channel, several media, no library}
\begin{enumerate}[leftmargin=1.6em,itemsep=2pt,topsep=2pt]
  \item Whether the wall yields, read as cells recovered against cells loaded.
  \item Messenger RNA molecules and genes per cell.
  \item Ribosomal fraction, which is the largest single lever on the cost of
    every later experiment and is a proportion, so it is answerable at any
    sequencing depth.
  \item Whether the media separate, which is both a biological result and the
    evidence that condition labels survive the protocol.
\end{enumerate}

\subsection{Three Things to Get Right at the Bench\secstatus{todo}}
\label{sec:counting}

\notesec{Count the cells on an instrument that can see yeast}
\org{S.\ cerevisiae} runs about 3.5 to 5\,\textmu m
across\external{vendor sizing notes; no primary in the mirror}, at or below the
stated floor of most brightfield counters, and 10x's own guidance is that cells
under 10\,\textmu m ``may be difficult to distinguish from debris'' and that a
fluorescent viability stain rather than trypan blue is ``strongly
recommended''\external{10x CG00053}. Two instruments are validated on this
organism by their makers, the Logos LUNA-FX7 and the ChemoMetec
NC-203\external{Logos and ChemoMetec application notes}. The Bio-Rad TC20
specifies a 6\,\textmu m floor\external{Bio-Rad TC20 brochure}, which excludes
yeast. The count sets the droplet loading rate and therefore the doublet rate,
and the core named it as the single most critical step.

\notesec{Replicate one genotype across two independent wells} Not a second
channel and not a technical resequencing. Two wells, or two cartridges, with
wild-type cells in each, so that a profile can be compared against itself with a
reference from its own batch. This is the check that would have caught the
published screen, it costs a handful of wells, and without it a null result
later cannot be told from a broken assay.

\notesec{Sequence twice, shallow then deep} A MiSeq run confirms the chemistry
worked, and the questions it answers are proportions, so they need almost no
depth. A NovaSeq X 10B split lane then buys 600 million read pairs for \$1{,}160,
which is less than the MiSeq run costs and 24 times the reads. There is no
earlier checkpoint than sequencing something: complementary DNA traces and
library concentration do not predict whether barcoding
worked~\citep{brettnerUltraHighthroughputMassively2024}.

\subsection{What Happens Next, Under Each Outcome\secstatus{todo}}
\label{sec:next}

\notesec{If $q$ is high in \org{S.\ cerevisiae}} The library needs no rebuilding.
Run a focused single-guide pilot on a metabolic sub-library at 100 cells per
gene, and add growth media to it through the plate route, where they are close
to free.

\notesec{If $q$ is near zero} The 5$'$ scaffold primer does not engage a
polymerase III transcript. Build the polyadenylated guide barcode and validate
it on about ten arrayed strains with known knockdown phenotypes against bulk RNA
sequencing of the same strains. Ground truth is only available while the design
is arrayed, since a pool cannot say what was in a given cell.

\notesec{If the \org{P.\ kudriavzevii} wall does not yield} The failure is in
sample preparation rather than in the platform, and the next attempt varies the
digestion rather than the sequencing. Nothing else in the plan depends on it,
which is why it is worth running now.

\notesec{If both work} The two tracks converge on one question, which is whether
to spend the next budget on more perturbations per cell or on more growth media.
The arithmetic favors media, and the model that would consume the data needs
both, so the first crossed design is a small metabolic panel across a handful of
conditions rather than a genome-scale screen in one.

\begin{keybox}
\textbf{Two runs, two people, about \$8{,}300, and every later decision depends
on what they return.} The \org{S.\ cerevisiae} channel decides whether a
construction project stands between us and a screen. The
\org{P.\ kudriavzevii} channel decides whether the assay works on that organism
at all. Build the three quality checks of \cref{sec:nadal} into both, because
the published screen they come from is the design this plan would otherwise have
copied.
\end{keybox}


\clearpage
\small
\bibliography{references}

\end{document}
